(1) 固定电荷时存储场能 \(U_Q=Q^2/(2C)=Q^2d/(2\varepsilon A)\)。对板距 \(d\) 的广义力定义为 \(F_d=-\partial U_Q/\partial d=-Q^2/(2\varepsilon A)<0\)，负号表示极板受力趋向缩小间距，即静电吸引。力的大小也等于电磁应力 \(\varepsilon E^2A/2\)，此时 \(E=Q/(\varepsilon A)\) 与 \(d\) 无关。此结果对断电孤立电容器成立，若固定电压而仍保持 \(Q\) 常数将违反控制条件。

(2) 恒压时 \(Q=C(d)V=\varepsilon AV/d\) 随 \(d\) 变动，场能 \(U_f=\frac12CV^2=\varepsilon AV^2/(2d)\)。令板距改变 \(dd\)，场能微分 \(dU_f=\frac12V^2dC\)，理想电源向电容器输送电功 \(dW_{\rm src}=VdQ=V^2dC\)。能量守恒 \(dU_f=dW_{\rm src}+dW_{\rm mech,on}\)，故外部机械手对电场所做功 \(dW_{\rm mech,on}=-\frac12V^2dC\)，相应电场在板距方向做功 \(dW_{\rm field}=\frac12V^2dC=-\varepsilon AV^2dd/(2d^2)\)。因此电场力 \(F_d=-\varepsilon AV^2/(2d^2)\) 仍为吸引力，而不是单纯取 \(-dU_f/dd>0\) 的错误结果。

(3) 将恒压电源纳入广义势能，适于固定电压的 Legendre 变换为 \(\mathcal F_V=U_f-VQ=-\frac12C(d)V^2=-\varepsilon AV^2/(2d)\)。其广义力为 \(F_d=-\partial\mathcal F_V/\partial d=+\frac12V^2(dC/dd)=-\varepsilon AV^2/(2d^2)\)，与能量收支独立一致。固定电压时必须使用该电势而非孤立电容器能量，因为电源可向或从电极输送电荷并携带功。

(4) 两种控制均有实际电场 \(E=Q/(\varepsilon A)=V/d\)，电场压力大小 \(p=\varepsilon E^2/2\)，总吸引力 \(F_d=-pA\)。若固定 \(Q\) 而缓慢拉开，所需外力正功进入逐渐增大的电场能。若固定 \(V\) 缓慢拉开，电容减少、流回电源的电荷为负 \(dQ<0\)，源收到能量 \(-dW_{\rm src}=-VdQ\)；外部正机械功与减少的场能一同供给电源，二者正好由前述微分式分配。
